\input{../tex-inputs/latex-leseansicht-vorspann}

\section[Gustav Schwarzkopf an Arthur Schnitzler, 31. 7. 1896]{L04287 Gustav Schwarzkopf an Arthur Schnitzler, 31. 7. 1896}
\nopagebreak\mylabel{L04287v}
\rehead{ }\normalsize\beginnumbering\briefempfaengerindex{Schnitzler, Arthur@\textsc{Schnitzler, Arthur}!zzzSchwarzkopf, Gustav@\emph{von Gustav Schwarzkopf}!1896-07-312@{31. 7. 1896}|(be}
\toendnotes[C]{\smallbreak\pagebreak[2]}
\correspDesc{Versand  durch Gustav Schwarzkopf am 31. 7. 1896 in Wien
\newline{}Zustellung  in Kopenhagen
\newline{}Erhalt  durch Arthur Schnitzler im Zeitraum [1. 8. 1896
                  – 5. 8. 1896?] in Kopenhagen}\toendnotes[C]{\smallbreak}
\Standort{CUL, Schnitzler, B 96a.}
\physDesc{Briefkarte, 1252 Zeichen
\newline{}Handschrift: schwarze Tinte, deutsche Kurrent}\toendnotes[C]{\smallbreak}
\pstart
           \raggedleft{}{\pb}31. VII 96\pend
           
\pstart{}Lieber Freund!\pend\vspace{0.5em}
\pstart
           Ich{ }ſoll heute Abend abreiſen. Mit \textsc{Chiavacci\pwindex{Chiavacci, Vincenz 15.\,6.\,1847 Wien – 2.\,2.\,1916 ebd.@\textsc{Chiavacci, Vincenz} (15.\,6.\,1847 Wien – 2.\,2.\,1916 ebd.), \emph{Schriftsteller, Journalist, Beamter}|pw}}. Ziel: Garda-See\oindex{Lago di Garda@\textbf{Lago di Garda}, \emph{See}|pw}, dann eventuell Stilfſerjoch\oindex{Stilfser Joch@\textbf{Stilfser Joch}, \emph{Berg}|pw} oder Mendelpaß\oindex{Mendelpass@\textbf{Mendelpass}, \emph{Pass}|pw}. Aber ich möchte noch nicht darauf{ }ſchwören, daß
               ich dieſen großen Plan auch wirklich durchführen werde. Nach Kopenhagen\oindex{Kopenhagen@\textbf{Kopenhagen}, \emph{Hauptstadt}|pw} wäre ich{ }ſehr gerne geko{\geminationm}en – das dürfen Sie glauben – aber \label{K_L04287-1v}\edtext{für dieſe Strecke gibt es für mich keinen
                  \textsc{Carlweiss\pwindex{Karlweis, Carl 23.\,11.\,1850 Wien – 27.\,10.\,1901 ebd.@\textsc{Karlweis, Carl} (23.\,11.\,1850 Wien – 27.\,10.\,1901 ebd.), \emph{Schriftsteller}|pw}}}{\lemma{\textnormal{\emph{für … Carlweiss}}}\Cendnote{\textnormal{Für die \emph{Südbahnstrecke}\orgindex{Südbahnstrecke@Südbahnstrecke|pwk} konnte Schwarzkopf\pwindex{Schwarzkopf, Gustav 7.\,11.\,1853 Wien – 13.\,11.\,1939 ebd.@\textsc{Schwarzkopf, Gustav} (7.\,11.\,1853 Wien – 13.\,11.\,1939 ebd.), \emph{Schriftsteller}|pwk}
                  über Karlweis\pwindex{Karlweis, Carl 23.\,11.\,1850 Wien – 27.\,10.\,1901 ebd.@\textsc{Karlweis, Carl} (23.\,11.\,1850 Wien – 27.\,10.\,1901 ebd.), \emph{Schriftsteller}|pwk} Freifahrkarten bekommen, vgl. XXXX Auszeichnungsfehler: Dokument L04280 nicht gefunden. }}}\label{K_L04287-1}. – Die
               Bevölkerungszahl hat indeſſen einigen Zuwachs erhalten, zwei Buben, einen Ludwig \textsc{Chiavacci}\pwindex{Chiavacci, Ludwig 9.\,7.\,1896 Wien – 14.\,8.\,1970 New York City@\textsc{Chiavacci, Ludwig} (9.\,7.\,1896 Wien – 14.\,8.\,1970 New York City), \emph{Psychiater, Dermatologe}|pw} und einen Franz \textsc{Hirschfeld}\pwindex{Hirschfeld, Franz *~27.\,7.\,1896@\textsc{Hirschfeld, Franz} (*~27.\,7.\,1896)|pw}, beide mit ziemlicher {\pb}Verſpätung
               eingelangt. – Was iſt’s mit \textsc{Lichtenberg\pwindex{Lichtenberg, Georg Christoph 1.\,7.\,1742 Ober-Ramstadt – 24.\,2.\,1799 Göttingen@\textsc{Lichtenberg, Georg Christoph} (1.\,7.\,1742 Ober-Ramstadt – 24.\,2.\,1799 Göttingen), \emph{Schriftsteller, Philosoph, Physiker}!Georg Christ. Lichtenbergs ausgewählte Schriften@\strich\emph{Georg Christ. Lichtenbergs ausgewählte Schriften}|pwv}\pwindex{Lichtenberg, Georg Christoph 1.\,7.\,1742 Ober-Ramstadt – 24.\,2.\,1799 Göttingen@\textsc{Lichtenberg, Georg Christoph} (1.\,7.\,1742 Ober-Ramstadt – 24.\,2.\,1799 Göttingen), \emph{Schriftsteller, Philosoph, Physiker}|pw}}? Sollte ich unbewußt die Höhe der? Sollte ich unbewußt die Höhe der »Moderne«
                  erklo{\geminationm}en und ein Plagiat begangen haben? Ich werde ja
               sehen, we{\geminationn} Sie mir freundlichſt die Stelle zeigen
               werden, de{\geminationn} das Buch\pwindex{Lichtenberg, Georg Christoph 1.\,7.\,1742 Ober-Ramstadt – 24.\,2.\,1799 Göttingen@\textsc{Lichtenberg, Georg Christoph} (1.\,7.\,1742 Ober-Ramstadt – 24.\,2.\,1799 Göttingen), \emph{Schriftsteller, Philosoph, Physiker}!Georg Christ. Lichtenbergs ausgewählte Schriften@\strich\emph{Georg Christ. Lichtenbergs ausgewählte Schriften}|pwv} auch noch kaufen, um mich davon zu überzeugen – nein.
               – Iſt \textsc{Richard B. H.\pwindex{Beer-Hofmann, Richard 11.\,7.\,1866 Wien – 26.\,9.\,1945 New York City@\textsc{Beer-Hofmann, Richard} (11.\,7.\,1866 Wien – 26.\,9.\,1945 New York City), \emph{Schriftsteller}|pw}} bei Ihnen? Grüßen Sie ihn herzlich von mir, ebenſo Dr \textsc{Goldma{\geminationn}\pwindex{Goldmann, Paul 31.\,1.\,1865 Breslau – 25.\,9.\,1935 Wien@\textsc{Goldmann, Paul} (31.\,1.\,1865 Breslau – 25.\,9.\,1935 Wien), \emph{Schriftsteller, Journalist}|pw}}. – Ihr Recitator, Marcell Salzer\pwindex{Salzer, Marcell 27.\,3.\,1873 Moravský Svätý Ján – 17.\,3.\,1930 Berlin@\textsc{Salzer, Marcell} (27.\,3.\,1873 Moravský Svätý Ján – 17.\,3.\,1930 Berlin), \emph{Schauspieler, Rezitator}|pw}, hat
               wieder für ihren Ruhm geſorgt, hat zweimal\eventindex{Pertl’s Drittes Kaffeehaus@\textbf{Pertl’s Drittes Kaffeehaus}!Lesung von Marcell Salzer, 17.7.1896@Lesung von Marcell Salzer, 17.7.1896|pwv}\eventindex{Hotel Zur Stadt Wien@\textbf{Hotel Zur Stadt Wien}!Lesung von Marcell Salzer, 5.8.1896@Lesung von Marcell Salzer, 5.8.1896|pwv} etwas von Ihnen geleſen (Denkſteine\pwindex{Schnitzler, Arthur 15.\,5.\,1862 Wien – 21.\,10.\,1931 ebd.@\textsc{Schnitzler, Arthur} (15.\,5.\,1862 Wien – 21.\,10.\,1931 ebd.), \emph{Schriftsteller, Mediziner}!Denksteine@\strich\emph{Denksteine}|pw}) einmal\eventindex{Pertl’s Drittes Kaffeehaus@\textbf{Pertl’s Drittes Kaffeehaus}!Lesung von Marcell Salzer, 17.7.1896@Lesung von Marcell Salzer, 17.7.1896|pwv} bei Pertl\oindex{Wien@\textbf{Wien}!II., Leopoldstadt@\textbf{II., Leopoldstadt}!Pertl’s Drittes Kaffeehaus@\textbf{Pertl’s Drittes Kaffeehaus}, \emph{Kaffeehaus}|pw} im
                  Prater\oindex{Wien@\textbf{Wien}!II., Leopoldstadt@\textbf{II., Leopoldstadt}!Prater@\textbf{Prater}, \emph{Park}|pw}, \label{K_L04287-2v}\edtext{einmal\eventindex{Hotel Zur Stadt Wien@\textbf{Hotel Zur Stadt Wien}!Lesung von Marcell Salzer, 5.8.1896@Lesung von Marcell Salzer, 5.8.1896|pwv} in Baden\oindex{Baden bei Wien@\textbf{Baden bei Wien}, \emph{Hauptstadt}|pw}}{\lemma{\textnormal{\emph{einmal in Baden}}}\Cendnote{\textnormal{Diese zweite Lesung\eventindex{Hotel Zur Stadt Wien@\textbf{Hotel Zur Stadt Wien}!Lesung von Marcell Salzer, 5.8.1896@Lesung von Marcell Salzer, 5.8.1896|pwkv} stand erst bevor, sie fand am
                     5. 8. 1896 im Hotel Zur Stadt Wien\oindex{Hotel Zur Stadt Wien@\textbf{Hotel Zur Stadt Wien}, \emph{Hotel}|pwk} statt.}}}\label{K_L04287-2}. Ihre Genossen waren H. So{\geminationn}enthal\pwindex{Sonnenthal, Hermine von 1.\,11.\,1863 Wien – 12.\,6.\,1922 ebd.@\textsc{Sonnenthal, Hermine von} (1.\,11.\,1863 Wien – 12.\,6.\,1922 ebd.), \emph{Schriftstellerin}!Mutterliebe@\strich\emph{Mutterliebe}|pwv}\pwindex{Sonnenthal, Hermine von 1.\,11.\,1863 Wien – 12.\,6.\,1922 ebd.@\textsc{Sonnenthal, Hermine von} (1.\,11.\,1863 Wien – 12.\,6.\,1922 ebd.), \emph{Schriftstellerin}|pw}, Anton Lindner\pwindex{Lindner, Anton 14.\,12.\,1874 Lviv – 30.\,12.\,1928 Wandsbek@\textsc{Lindner, Anton} (14.\,12.\,1874 Lviv – 30.\,12.\,1928 Wandsbek), \emph{Schriftsteller}!Fern von Wien@\strich\emph{Fern von Wien}|pwv}\pwindex{Lindner, Anton 14.\,12.\,1874 Lviv – 30.\,12.\,1928 Wandsbek@\textsc{Lindner, Anton} (14.\,12.\,1874 Lviv – 30.\,12.\,1928 Wandsbek), \emph{Schriftsteller}|pw} u. – Chiavacci\pwindex{Chiavacci, Vincenz 15.\,6.\,1847 Wien – 2.\,2.\,1916 ebd.@\textsc{Chiavacci, Vincenz} (15.\,6.\,1847 Wien – 2.\,2.\,1916 ebd.), \emph{Schriftsteller, Journalist, Beamter}!Rütli am Thury@\strich\emph{Das Rütli am Thury}|pwv}\pwindex{Chiavacci, Vincenz 15.\,6.\,1847 Wien – 2.\,2.\,1916 ebd.@\textsc{Chiavacci, Vincenz} (15.\,6.\,1847 Wien – 2.\,2.\,1916 ebd.), \emph{Schriftsteller, Journalist, Beamter}|pw} – auch als Junger angekündigt.–\pend
           
\pstart
           Ein neuer großer Dichter iſt in den letzten 14 Tagen nicht entdeckt worden. Die Zeit\orgindex{Zeit. Wiener Wochenschrift@Die Zeit. Wiener Wochenschrift|pw}{ }ſchläft den So{\geminationm}erſchlaf.\pend
           
\pstart
           Herzlichſt Ihr{\\[\baselineskip]}\spacefill\mbox{Gustav Schwk.}\pend
           \leftskip=0em{}\selectlanguage{ngerman}\endnumbering\briefempfaengerindex{Schnitzler, Arthur@\textsc{Schnitzler, Arthur}!zzzSchwarzkopf, Gustav@\emph{von Gustav Schwarzkopf}!1896-07-312@{31. 7. 1896}|)be}\mylabel{L04287h}
\begin{anhang}
\end{anhang}\newcommand{\dateiname}{L04287}\newcommand{\titel}{Gustav Schwarzkopf an Arthur Schnitzler, 31. 7. 1896}\newcommand{\editorInnen}{Herausgegeben von Jahnke, SelmaMüller, Martin Anton}\input{../tex-inputs/latex-leseansicht-abspann}
